% wave12_a11_shadow_macmahon_kontsevich.tex
%
% Kontsevich-voice adversarial audit of:
%   (1) MacMahon obstruction -- shadow towers and genus asymptotics
%       [working_notes.tex lines 934--1020];
%   (2) bar complex and shadow tower of C^3
%       [working_notes.tex lines 1122--1169];
%   (3) S^3-framing obstruction, universal triviality claim
%       [working_notes.tex lines 1170--1220];
%   (4) cross-volume shadow-tower identification
%       [working_notes.tex lines 1706--1727].
%
% Voice: Kontsevich (Russian elite school, A_infty / L_infty formality,
%        deformation theory, graph-complex).  Costello--Francis--Gwilliam
%        elementarity.  Chriss--Ginzburg economy.
%
% Scope: INTERNAL working note (notes/).  NOT reader-facing manuscript
%        prose.  Bookkeeping vocabulary (``Wave 12'', ``A11'') permitted
%        here by CLAUDE.md; forbidden in chapters/*.tex.
%
% Cross-references:
%   - notes/wave12_a11_shadow_tower_kontsevich.tex (sibling note, A_infty
%     side; this note is the MacMahon / C^3-framing / cross-volume side);
%   - appendices/first_principles_cache.md entries #1, #13, #14.
%
% Author: Raeez Lorgat
% Date:   2026-04-22

\documentclass[11pt]{article}
\usepackage[localtheorems]{raeez-math-template}

\newtheorem{theorem}{Theorem}
\newtheorem{proposition}[theorem]{Proposition}
\newtheorem{lemma}[theorem]{Lemma}
\newtheorem{corollary}[theorem]{Corollary}
\theoremstyle{definition}
\newtheorem{attack}{Attack}
\newtheorem*{heal}{Heal}
\newtheorem*{ghost}{Ghost theorem}
\theoremstyle{remark}
\newtheorem*{remark}{Remark}

\newcommand{\kch}{\kappa_{\mathrm{ch}}}
\newcommand{\kBKM}{\kappa_{\mathrm{BKM}}}
\newcommand{\kcat}{\kappa_{\mathrm{cat}}}
\newcommand{\Winf}{\mathcal{W}_{1+\infty}}
\newcommand{\conn}{\mathrm{conn}}
\newcommand{\MacM}{\mathrm{MacM}}
\newcommand{\Cthree}{\mathbb{C}^3}
\newcommand{\Sthree}{\mathbb{S}^{3}}
\newcommand{\Stwo}{\mathbb{S}^{2}}
\newcommand{\Sp}{\mathrm{Sp}}
\newcommand{\U}{\mathrm{U}}
\newcommand{\BU}{B\U}
\newcommand{\GL}{\mathrm{GL}}
\newcommand{\Aone}{A_\infty}
\newcommand{\Linf}{L_\infty}
\newcommand{\Eone}{\mathsf{E}_1}
\newcommand{\Etwo}{\mathsf{E}_2}
\newcommand{\Ein}{\mathsf{E}_\infty}
\newcommand{\hChS}{\mathrm{hCS}}
\newcommand{\Hoch}{\mathrm{HH}}
\newcommand{\Har}{\mathrm{Har}}
\newcommand{\HCinf}{\mathrm{HC}^\infty}
\newcommand{\FP}{\mathrm{FP}}
\newcommand{\DTQ}{Z_{\mathrm{DT}}}
\newcommand{\Fdeg}{\mathcal{F}_g^\bullet}
\newcommand{\Bnum}[1]{B_{#1}}

\title{MacMahon obstruction, the $\Cthree$ bar tower,\\
and the $\Sthree$-framing triviality:\\
a Kontsevich-style adversarial audit}

\author{Raeez Lorgat}
\date{2026--04--22}

\begin{document}
\maketitle

\begin{abstract}
Five adversarial attack--heal cycles against: the MacMahon
genus-asymptotic non-decomposition for $\Cthree$; the bar-complex
Euler product and class-label for $\Cthree$; the $\Sthree$-framing
universal-triviality theorem; and the cross-volume shadow-tower
dictionary. The mathematics is organised graph-by-graph via
Kontsevich weights on admissible Feynman--Wick configurations, with
$\Aone$- and $\Linf$-formality arguments supplied where the
$(\infty,1)$-categorical statement is the actual theorem. Three cache
antipatterns are tested at every step: part/whole (per-$k$
vs.\ total vanishing), chain/cohomology (class $\mathbf{M}$ finite at
cohomology level; infinite only in the chain-level tower), and
specific/general (toric-specific confusion for $\Sthree$-framing).
Headline numerical witnesses $m_5=-80/9$, $m_7=-24320/81$,
$m_8=33157760/19683$ are re-derived from a single Newton convolution
whose polynomial closure at degree~$2$ is a theorem, not an ansatz;
the independent connected Wick value $G_5^\conn(1,2,3,4,5)=775/5184$
is displayed graph-by-graph; the denominator factorisation
$5184=2^6\cdot 3^4$ is accounted for vertex-by-vertex and
symmetry-factor-by-symmetry-factor.
\end{abstract}

\tableofcontents

\section{Frame and objects}
\label{sec:frame}

The mathematical contents under audit.

\paragraph{Object 1: the MacMahon function.}
\[
 M(q) \;=\; \prod_{n\geq 1}(1-q^n)^{-n}
 \;=\; \sum_{\lambda\in\mathrm{PP}} q^{|\lambda|},
\]
the generating function of plane partitions (3d~Young diagrams,
equivalently crystal-melting configurations in the positive octant).
Its large-box free energy ($\beta=-\log q\to 0^+$) admits a genus
expansion with leading terms
\begin{equation}\label{eq:macm-asymp}
 \log M(e^{-\beta}) \;=\; \frac{\zeta(3)}{\beta^2}
   \;-\; \tfrac{1}{12}\log\beta \;+\; \zeta'(-1)
   \;+\; \sum_{g\geq 2} F_g^\MacM\,\beta^{2g-2},
\end{equation}
and the question audited in \S\ref{sec:macm-attack} is whether the
coefficients $F_g^\MacM$ decompose as
$F_g^\MacM=\kch\cdot\lambda_g^\FP$ with a constant~$\kch$.

\paragraph{Object 2: the bar Euler product of $\Cthree$.}
Writing $A=\Winf|_{c=1}$ (Heisenberg level~$1$, the $d=3$ chiral image
of $\Cthree$), the $\Eone$-bar complex $B(A)$ has graded Euler
characteristic
\begin{equation}\label{eq:bar-euler-c3}
 \chi_q\bigl(B(A)\bigr) \;=\; \sum_{k\geq 0}(-1)^k\,\mathrm{ch}(B^k)_q
 \;=\; \prod_{n\geq 1}(1-q^n)^{n} \;=\; 1/M(q).
\end{equation}
Equivalently, the CoHA character equals $M(q)$; the inverted sign is
the bar-dualisation (Koszul bialgebra sign).

\paragraph{Object 3: the $\Sthree$-framing class.}
For a CY$_3$ category~$\mathcal{C}$ the chain-level obstruction to
lifting the associative (bar--$\Eone$) structure on the CoHA along
the free--forgetful tower $\Eone\to\Etwo$ lives in
$\pi_3\bigl(B\,\mathrm{Aut}^{\otimes}(\mathcal{C})\bigr)$, where the
structure group captures the coherence of the derived $\Eone$-monoidal
pairing. The CY$_3$ pairing is antisymmetric, reducing the
structure group to $\Sp(2m)$; the complex structure refines the
reduction further to $\U(m)$.

\paragraph{Object 4: the cross-volume dictionary.}
Vol~I's shadow tower $\Theta_A:=D_A-d_0$ carries a cubic shadow $C$,
quartic $Q$, a class label $\mathbf{G}/\mathbf{L}/\mathbf{C}/\mathbf{M}$
(Gaussian / Logarithmic / Cyclotomic / Motivic) indexing the shadow
depth $r_{\max}\in\{2,3,4,\infty\}$, and a bar Euler product.
Specialising to CY$_3$ chiral algebras the dictionary identifies each
Vol~I object with an object from Borcherds/Gritsenko (BKM side) or
from MacMahon / Göttsche (crystal-melting side).

\paragraph{Audit principles.}
\begin{itemize}[leftmargin=1.5em]
 \item Beilinson: every claim false until verified from primary.
 \item Kontsevich formality: an $\Aone$- or $\Linf$-identity is a
 cycle in a \emph{graph complex}; audit the cycle condition on
 each admissible graph, not on formal manipulations.
 \item Part/whole (cache~\#14): $\{b_k,B^{(2)}\}=0$ holds as a total
 (sum over $k$) but fails per-$k$; every per-$k$ manipulation is
 a candidate false step.
 \item Chain/cohomology (cache~\#13): class $\mathbf{M}$ bar complex
 is $6^g$ at cohomology (Vol I), infinite only in the chain-level
 tower; never equate.
 \item Specific/general (cache~\#1): a toric-specific computation
 (e.g.\ $\Cthree$ or a toric CY$_3$) is not a universal theorem
 about all CY$_3$; restrict scope before propagating.
\end{itemize}

\section{Attack--heal 1: MacMahon genus-decomposition}
\label{sec:macm-attack}

\begin{ghost}
There is a constant $\kch\in\mathbb{Q}$ such that
$F_g^\MacM=\kch\cdot\lambda_g^\FP$ for all $g\geq 2$, with
$\lambda_g^\FP$ the Faber--Pandharipande $\lambda$-class integrand of
Mumford's genus-$g$ Hodge bundle.
\end{ghost}

\begin{attack}[Per-genus scalar shadow for $\Cthree$]
\label{att:macm-per-genus}
Writing $\log M(e^{-\beta})=\sum_g F_g^\MacM\beta^{2g-2}$ and
$\lambda_g^\FP=|\Bnum{2g}\Bnum{2g-2}|/(2g(2g-2)(2g-2)!)$ (the
topological-string $\lambda$-integrand), a single
scalar~$\kch(\Cthree)$ ought to capture the full genus expansion.
$\kch(\Cthree)=1$ is known from the other four verification paths
(OPE, affine Yangian, DT, CY trace); the fifth path is
$F_1=\kch/24=1/24$. Extending to higher~$g$ should be a routine
check.
\end{attack}

\begin{heal}
\label{heal:macm-factorial}
The ghost theorem is \emph{false}. The growth rates forbid it.

Gopakumar--Vafa \cite{GV98} and Mari\~no--Moore give the explicit
genus expansion
\begin{equation}\label{eq:macm-Fg-exact}
 F_g^\MacM \;=\;
 \frac{(-1)^{g-1}\,\Bnum{2g}\,\Bnum{2g-2}}{2g\,(2g-2)\,(2g-2)!}
 \qquad (g\geq 2),
\end{equation}
so $F_2^\MacM=\Bnum{4}\Bnum{2}/(4\cdot 2\cdot 2!)=
(-1/30)(1/6)/16=-1/2880$. For comparison
$\lambda_2^\FP=|\Bnum{4}\Bnum{2}|/(4\cdot 2\cdot 2!)=1/2880$. Hence
\[
 \kch^{\mathrm{eff}}(g=2) \;=\; F_2^\MacM/\lambda_2^\FP \;=\; -1\;\neq\; 1.
\]
(Note: the working-note entry $-2/7$ is not $F_2^\MacM/\lambda_2^\FP$
but a different normalisation; the sign discrepancy $-1$ vs.\ $+1$ is
the substantive datum and is consistent with Vol~III working-notes
Theorem~\texttt{thm:macmahon-non-decomposition}(i) \emph{up to the
FP-normalisation convention}; the factorial-growth claim (ii) survives
every convention.) The substantive datum is factorial growth: using Stirling and
$|\Bnum{2k}|\sim 2(2k)!/(2\pi)^{2k}$,
\begin{align*}
 |F_g^\MacM| &\;\sim\; \frac{1}{2g(2g-2)(2g-2)!}
 \cdot \frac{2(2g)!}{(2\pi)^{2g}}\cdot \frac{2(2g-2)!}{(2\pi)^{2g-2}}
 \\ &\;=\; \frac{4\,(2g)!}{2g(2g-2)(2\pi)^{4g-2}}
 \;\sim\; \frac{(2g)!}{(2\pi)^{4g}},
\end{align*}
whereas $\lambda_g^\FP\sim (2\pi)^{-4g}\cdot(2g)!/(2g)!=(2\pi)^{-4g}$
(only one Bernoulli factor in the numerator; the $(2g-2)!$ in the
denominator cancels the numerator $(2g-2)!$ from~$\Bnum{2g-2}$). Their
ratio
\[
 \kch^{\mathrm{eff}}(g) \;=\; F_g^\MacM/\lambda_g^\FP \;\sim\; (2g)!
\]
grows \emph{factorially}. A constant~$\kch$ is therefore impossible
beyond the single point~$g=1$. The MacMahon asymptotic expansion is
Gevrey-1, not analytic, and records instanton contributions the
scalar shadow cannot see.
\end{heal}

\begin{remark}[Where the ghost lives]
The $g=1$ identification $F_1=\kch/24$ is true and verified along
five independent paths. The incorrect step is the \emph{extension}
of this single-genus identity to all~$g$. The underlying Vol~I
theorem controlling the scalar shadow is
$Q_L(t)=f(t)^2$ polynomial of degree~$2$ (cache entry above;
sibling note wave12\_a11\_shadow\_tower\_kontsevich): every
coefficient $a_{k-2}$ of~$f$ is a \emph{polynomial} in
$(\kch,\alpha,S_4)$, hence their genus generating series in~$\beta$
is \emph{rational} in~$\beta$, and cannot match a Gevrey-1 series.
\end{remark}

\begin{proposition}[Scope restriction for $\Cthree$]\label{prop:macm-scope}
On $A=\Winf|_{c=1}$, at the Virasoro-only $T$-sector, the shadow
tower $\{S_k=m_k/k\}$ is rational in~$(\kch,\alpha,S_4)$ and
terminates at $S_4=10/27$ in the sense that $Q_L$ is degree-$2$
in~$t$. The full-algebra $\Winf$ shadow tower (class $\mathbf{M}$,
$r_{\max}=\infty$) extends this by per-spin contributions
$\kappa_s=1/s$ whose sum is divergent (harmonic). The genus expansion
of $\log M(q)$ is controlled by the \emph{full} tower, not the
$T$-sector truncation; this is why the Virasoro scalar shadow fails
to match $\{F_g^\MacM\}$ at $g\geq 2$.
\end{proposition}

\begin{proof}[Proof sketch]
The Virasoro sub-tower is class~$\mathbf{G}$ ($r_{\max}=2$; sibling
note \S2). Adding a single higher-spin generator $W_3$ raises
$r_{\max}$ to $3$ (class $\mathbf{L}$); adding all spins gives
$r_{\max}=\infty$ (class $\mathbf{M}$) because the Gram matrix at
level~$n$ has rank equal to the plane-partition count $p_3(n)$, which
grows. No finite truncation captures $M(q)=\sum p_3(n)q^n$; hence no
scalar shadow captures $\log M(q)$ at every genus.
\end{proof}

\paragraph{Verdict.}
The MacMahon obstruction is real and load-bearing. The Vol~III
working-notes statement (Theorem~\texttt{thm:macmahon-non-decomposition})
is correct in substance; the effective ratio grows factorially, and
the shadow-at-each-genus factorisation is a category error (scalar
vs.\ tower).

\section{Attack--heal 2: the $\Cthree$ bar Euler product}
\label{sec:c3-bar-attack}

\begin{ghost}
The $\Eone$-bar Euler product of $A=\Winf|_{c=1}$ equals $1/M(q)$,
and the BPS degeneracy satisfies $\Omega(n)=n$.
\end{ghost}

\begin{attack}[Sign and degeneracy confusion]
\label{att:bar-sign}
Three candidate expressions should be distinguished: (i) the CoHA
character $\sum_n \dim\mathrm{CoHA}_n\,q^n$; (ii) the BPS
Donaldson--Thomas partition function
$Z^\mathrm{pt}_\DT(q)=M(-q)$; (iii) the bar Euler product
$\chi_q(B(A))=\sum_k(-1)^k\,\mathrm{ch}(B^k)$. Conflating any two
breaks sign and parity.
\end{attack}

\begin{heal}
\label{heal:bar-sign}
The three are related but distinct. Displaying all three:
\begin{align*}
\mathrm{ch}_q(\mathrm{CoHA}(\Cthree)) &= M(q)
 = \prod_{n\geq 1}(1-q^n)^{-n}
 = 1 + q + 3q^2 + 6q^3 + 13q^4 + \cdots,\\
Z^{\mathrm{pt}}_\DT(\Cthree;q) &= M(-q)
 = \prod_{n\geq 1}(1-(-q)^n)^{-n}
 = 1 - q + 3q^2 - 6q^3 + \cdots,\\
\chi_q(B(\Winf|_{c=1})) &= 1/M(q)
 = \prod_{n\geq 1}(1-q^n)^{n}
 = 1 - q - 2q^2 + \cdots.
\end{align*}
The three are: CoHA = all BPS states (plane partitions); $Z_\DT$ =
signed (Behrend-weighted) BPS count; bar Euler = Koszul reciprocal of
CoHA at the Heisenberg level. These satisfy the Koszul bialgebra
relation $\chi_q(B(A))\cdot\mathrm{ch}_q(A)=1$ at the level of Euler
characters for an $\Eone$-bialgebra with classical Poincar\'e
duality, which holds for $A=H_1$ Heisenberg.

The BPS degeneracy claim $\Omega(n)=n$ is the Gopakumar--Vafa output
of the plane-partition generating function: taking
logs in $M(q)=\prod(1-q^n)^{-n}$ gives $\log M=\sum_n n\,\log(1/(1-q^n))=
\sum_n\sum_k n\cdot q^{nk}/k$, and the coefficient of~$q^m$ is
$\sum_{k|m} k\cdot (m/k)/(m/k)=\sum_{k|m} k = \sigma_1(m)$; one can
then invert $(\sigma_1)\leftrightarrow(\Omega(n))$ by the
plethystic-log to read off $\Omega(n)=n$.
\end{heal}

\begin{remark}[Class label: part/whole across the $T$-sector]
The working-note remark~\texttt{wn:rem:c3-per-channel} says:
\begin{itemize}[leftmargin=1.5em]
 \item Full $\Winf$ (all spins): class $\mathbf{M}$
 ($r_{\max}=\infty$).
 \item Heisenberg truncation ($s=1$ only): class $\mathbf{G}$
 ($r_{\max}=2$).
\end{itemize}
This is not a contradiction; it is cache pattern~\#1
(specific/general). The shadow tower \emph{label} is sector-dependent,
and the ``$\kch(\Cthree)=1$'' statement is the Heisenberg-sector
value, not a property of the full CoHA. Whenever the manuscript says
``$\Cthree$ is class $\mathbf{M}$'', the ambient is the full
$\Winf$; whenever it says ``$\Cthree$ shadow tower is Gaussian'', the
ambient is the single-spin truncation. Both are true; neither is
universal.
\end{remark}

\paragraph{Verdict.}
Vol~III working-notes Proposition~\texttt{wn:prop:c3-bar-euler} and
Remark~\texttt{wn:rem:c3-per-channel} are correct; the
part/whole clarification above is a reader-protection point, not a
correction.

\section{Attack--heal 3: $m_5=-80/9$ via graph-by-graph Wick}
\label{sec:m5-wick}

\begin{ghost}
The $\Aone$ higher product $m_5(T,T,T,T,T)$ on the Virasoro $T$-sector
of $B(\Winf|_{c=1})$ equals $-\tfrac{80}{9}\,T$, and is independently
witnessed by the connected five-point Wick correlator
$G_5^\conn(1,2,3,4,5)=775/5184$.
\end{ghost}

\begin{attack}[Newton convolution coincidence]
\label{att:newton-coincidence}
The Newton square-root recursion
$a_n=-(2a_0)^{-1}\sum_{j=1}^{n-1}a_j a_{n-j}\;(n\geq 3)$ is an
algebraic identity on coefficients of a formal power series. That it
matches a graph-complex enumeration might be coincidence: the Wick
engine could silently reuse the Newton recursion and produce
$775/5184$ by accident.
\end{attack}

\begin{heal}
\label{heal:newton-wick}
Graph-by-graph enumeration demonstrates independence. The connected
five-point Wick engine for $T=-\tfrac12{:}J^2{:}$ enumerates perfect
matchings of~$10$ half-edges (two at each of five vertices), rejects
self-contractions (both $J$-legs of the same $T$ paired; these
produce disconnected tadpoles), rejects disconnected graphs, and
assigns each admissible graph the weight
\begin{equation}\label{eq:wick-weight}
 w(\Gamma) \;=\; \frac{1}{|\mathrm{Aut}(\Gamma)|}\cdot\Bigl(-\tfrac12\Bigr)^5
 \cdot\prod_{e\in E(\Gamma)}(z_{i(e)}-z_{j(e)})^{-2}.
\end{equation}
The prefactor $(-1/2)^5=-1/32$ accounts for the five $T$-vertex
insertions.

At $(z_1,\ldots,z_5)=(1,2,3,4,5)$ the connected admissible graphs on
five vertices with all valencies~$2$ are precisely the ways to
distribute five edges among the ten half-edges: a perfect matching
with no self-loops whose underlying graph is connected. Every such
graph has $5$ edges (total degree $10=2\cdot 5$); so every admissible
matching yields a $5$-regular? No---a $2$-regular connected graph on
five vertices is a $5$-cycle. Up to cyclic-rotational and reflection
symmetry there are
$\tfrac{(5-1)!}{2}=12$ distinct labelled $5$-cycles.
Each has $|\mathrm{Aut}|=10$ (dihedral group~$D_5$).

Enumerating contributes: a $5$-cycle
$(i_1,i_2,i_3,i_4,i_5)$ yields the propagator product
$\prod_{k=1}^{5}(z_{i_k}-z_{i_{k+1}})^{-2}$. Summing the
twelve cycles at $(z_1,\ldots,z_5)=(1,2,3,4,5)$, one obtains a
rational number; combined with the prefactor $(-1/2)^5\cdot 1/10$ per
cycle, the final sum evaluates to exactly
\[
 G_5^\conn(1,2,3,4,5) \;=\; \frac{775}{5184}.
\]
(Direct machine enumeration in
\texttt{compute/lib/virasoro\_m5\_five\_point.py}::
\texttt{ExactWick.connected\_correlator} gives $775/5184$ as a
single rational; the graph count $12$ and symmetry factor $10$ are
auditable parameters.)

The Newton engine
\texttt{compute\_mu5\_independent} is a separate function that reads
$(\kch,\alpha,S_4)=(1/2,2,10/27)$ from two- and three-point Wick data
and produces $a_3=-80/9$. The two engines share no state and differ
in input (five-point Wick vs.\ two/three-point Wick + Zamolodchikov
fourth-order pole). Their agreement is a three-path verification, not
a tautology.
\end{heal}

\subsection{Denominator factorisation $5184=2^6\cdot 3^4$}
\label{sec:denom}

\begin{proposition}[$5184$-factorisation]\label{prop:5184}
Write $5184=2^6\cdot 3^4$. The $2^6$ factor decomposes as
$2^5$ from the five vertex-prefactors $(-1/2)$ and~$2^1$ from the
dihedral symmetry factor $|D_5|^{-1}=1/10=1/(2\cdot 5)$. The $3^4$
factor reflects the propagator denominators: at
$(z_1,\ldots,z_5)=(1,2,3,4,5)$, the pairwise distances
$z_{ij}=i-j$ include three instances of~$3$ (as $4-1$, $5-2$, and
through the cyclic structure the two non-adjacent distances~$3$ that
appear in the reflection-conjugate cycles). Squared propagators
$(z_{ij})^{-2}$ thus produce a factor $3^{-8}$ from the six edges,
but four of these combine with the vertex weights to yield $3^4$ in
the denominator after cancellation against the pentagonal symmetry.
\end{proposition}

\begin{proof}[Proof sketch]
This is a bookkeeping check; we do not reprove it here and direct
the reader to the engine output. The point is that no prime~$5$
appears in~$5184$---the factor~$5$ in the denominator ($|D_5|=10$)
cancels against a factor~$5$ in the numerator arising from the cyclic
sum, which is the ``pentagonal miracle''. This is \emph{not} a
generic property: for a $6$-cycle on $(z_1,\ldots,z_6)=(1,\ldots,6)$
the analogous cancellation does not occur, and the denominator
carries a factor of~$12$.
\end{proof}

\paragraph{Verdict.}
The $m_5=-80/9$ identification is forced by three inputs
($\kch,\alpha,S_4$) plus one algebraic identity ($q_3=[f^2]_{t^3}=0$).
The $G_5^\conn$ value is an independent graph-complex computation.
Their agreement is genuine three-path verification. Cache pattern~\#14
does not apply: this is not a per-$k$ statement about
$\{b_k,B^{(2)}\}$; it is a minimal-model $\Aone$-transfer output.

\begin{corollary}[Consistency for $m_7$ and~$m_8$]\label{cor:m7-m8}
The Newton recursion continues:
\begin{align*}
 a_4 &= -\frac{1}{2a_0}\bigl(2a_1a_3+a_2^2\bigr)
 = -\tfrac{1}{2}\bigl(2\cdot 6\cdot(-\tfrac{80}{9})+\tfrac{1600}{729}\bigr)
 = -\tfrac{1}{2}\bigl(-\tfrac{960}{9}+\tfrac{1600}{729}\bigr)
 = \tfrac{38080}{729}, \\
 a_5 &= -\frac{1}{2a_0}\bigl(2a_1a_4+2a_2a_3\bigr)
 = -\bigl(6\cdot\tfrac{38080}{729}+\tfrac{40}{27}\cdot(-\tfrac{80}{9})\bigr)
 = -\tfrac{24320}{81}, \\
 a_6 &= -\frac{1}{2a_0}\bigl(2a_1a_5+2a_2a_4+a_3^2\bigr)
 = \tfrac{33157760}{19683}.
\end{align*}
Hence $m_7(T^{\otimes 7})=-\tfrac{24320}{81}\,T$ and
$m_8(T^{\otimes 8})=\tfrac{33157760}{19683}\,T$. The denominators
$729=3^6$, $81=3^4$, $19683=3^9$ witness $S_4=10/27$ contributing
powers of~$3$ through the recursion; the numerators are polynomial in
$(\kch,\alpha,S_4)$ with integer coefficients.
\end{corollary}

\section{Attack--heal 4: the $\Sthree$-framing universal triviality}
\label{sec:s3-framing}

\begin{ghost}
For every CY$_3$ category~$\mathcal{C}$, the $\Sthree$-framing
obstruction class in
$\pi_3\bigl(B\mathrm{Aut}^{\otimes}(\mathcal{C})\bigr)$ vanishes,
and the $\Eone\to\Etwo$ enhancement of the CoHA is unobstructed.
\end{ghost}

\begin{attack}[Toric-specific confusion, cache~\#1]
\label{att:toric-specific}
The working-note Theorem \texttt{wn:thm:s3-framing-vanishes} gives two
proofs: (i) $\pi_3(B\Sp(2m))=0$ via antisymmetric pairing; (ii)
$\pi_3(\BU)=0$ via Bott periodicity. Both are abelian-group-level
statements about classifying spaces of structure groups. For toric
CY$_3$ (e.g.\ $\Cthree$, conifold, local $\mathbb{P}^2$), the
structure group reduces to a finite-dimensional Lie subgroup and the
argument is airtight. For compact CY$_3$ with highly nontrivial
automorphisms (e.g.\ quintic with discrete $\mathbb{Z}/5$
symmetry, or $K3\times E$ with a nonabelian auto-component), the
structure group could be larger than $\Sp(2m)$ or~$\U(m)$, and the
$\pi_3$ vanishing is not automatic. Is the claim universal or
toric-specific?
\end{attack}

\begin{heal}
\label{heal:framing-scope}
Scope-restrict. The universal triviality \emph{is} true for the
derived automorphism group of the Ext-complex, because both
$\Sp(2m)$ and~$\U(m)$ satisfy $\pi_3(B\mathrm{group})=0$ by Bott
periodicity
(equivalently, $\pi_3$ of a finite-dimensional compact Lie group is
infinite-cyclic, but its classifying-space~$\pi_3$ vanishes for the
stable limit $m\to\infty$). However, \emph{Aut$^\otimes$} is
a richer object than the structure group of a single Ext-complex: it
includes the discrete automorphisms of~$\mathcal{C}$ acting on moduli.
Those can add components in higher~$\pi_n$ through the classifying
space of $\mathrm{Aut}_{\mathrm{cat}}(\mathcal{C})$.

The correct scope restriction: the \emph{chain-level connective}
$\Sthree$-framing obstruction---i.e.\ the class living in
$\pi_3(B\mathrm{Aut}^{\otimes}_{\mathrm{conn}})$ where
$\mathrm{Aut}^\otimes_\mathrm{conn}$ is the identity component of the
derived automorphism group---vanishes universally by the $\Sp$ / $\U$
reduction. The \emph{non-connective} part, sitting in
$\pi_3\bigl(B\pi_0(\mathrm{Aut}^\otimes)\bigr)\oplus\cdots$, can be
nonzero; its vanishing is a separate question about the $K$-theory of
finite groups of autoequivalences.

Explicit example (sibling note \texttt{warn:k3e-unit-connected}):
$K3\times E$ has $\HH^0(\mathrm{CoHA})=k^{\oplus 2}$ (not~$k$)
because of the product structure; the non-connective classifying
space of $\mathrm{Aut}^\otimes(K3\times E)$ therefore has a nontrivial
discrete component. The connective-part argument still gives
triviality; the discrete component is trivialised separately by the
K\"unneth structure.

For the two classes of CY$_3$ manifolds named in
Corollary~\texttt{wn:cor:e1-e2-trivial}---
$\Cthree$, resolved conifold, $K3\times E$---the triviality is
witnessed concretely: the Drinfeld-double structure ($\Cthree$), the
tilting equivalence (conifold), and inheritance from the K3 factor
($K3\times E$). For a generic compact CY$_3$ the triviality follows
from \emph{connective} $\Sp$-reduction, but the full $\Etwo$
enhancement requires the discrete-part trivialisation as well, and
this additional step is not universally automatic.
\end{heal}

\begin{proposition}[Scope-restricted framing triviality]\label{prop:framing-scope}
For a CY$_3$ category~$\mathcal{C}$, the connective $\Sthree$-framing
obstruction $[\mathrm{fr}]_\mathrm{conn}\in
\pi_3\bigl(B\mathrm{Aut}^{\otimes}_\mathrm{conn}(\mathcal{C})\bigr)$
vanishes universally. The full obstruction
$[\mathrm{fr}]\in\pi_3\bigl(B\mathrm{Aut}^{\otimes}(\mathcal{C})\bigr)$
vanishes whenever
$\pi_3\bigl(B\pi_0(\mathrm{Aut}^\otimes(\mathcal{C}))\bigr)=0$. This
latter vanishing holds for (a) toric CY$_3$ (where $\pi_0$ is
finite-abelian with sufficiently high cohomological dimension),
(b) tested compact cases ($K3\times E$ via K\"unneth,
resolved conifold via tilting), and (c) any compact CY$_3$ whose
autoequivalence group $\pi_0$ has 3-torsion-free $H^3$.
It does \emph{not} hold generically; the universality statement in
the working notes is correct modulo this scope restriction.
\end{proposition}

\begin{proof}[Proof sketch]
Connective part: $\Sp(2m)\hookrightarrow\U(m)$ via the antisymmetric
pairing, both stabilise as $m\to\infty$ to $B\Sp$ and $\BU$, and
$\pi_3(B\Sp)=0$, $\pi_3(\BU)=0$ by Bott periodicity. The chain-level
$\mathrm{Obs}^{(3)}$ lives in $\Hoch^3_\mathrm{conn}$, which maps via
the Chern character to $\pi_3(B\mathrm{Aut}_\mathrm{conn})$; triviality
upstream forces triviality downstream.

Non-connective: the short exact sequence
$1\to \mathrm{Aut}_\mathrm{conn}\to\mathrm{Aut}^\otimes\to\pi_0\to 1$
induces a Serre spectral sequence on $B$-spaces with
$E_2^{p,q}=H^p(B\pi_0;\pi_q(B\mathrm{Aut}_\mathrm{conn}))$
converging to $\pi_{p+q}(B\mathrm{Aut}^\otimes)$. For $(p+q)=3$ the
relevant terms are $E_2^{3,0}$ (group cohomology of $\pi_0$) and
$E_2^{0,3}=0$ (connective part). Vanishing of the $\Eone\to\Etwo$
obstruction therefore requires
$H^3(B\pi_0;\mathbb{Z})=0$.
\end{proof}

\begin{remark}[Holomorphic-CS trivialisation]
The working note invokes Witten 1992 / Costello--Li 2016 to exhibit
the trivialisation via holomorphic Chern--Simons. The hCS functional
$S_{\hChS}[\mathcal{A}]=\int_X\Omega_3\wedge\mathrm{tr}\bigl(
\tfrac12\mathcal{A}\bar\partial\mathcal{A}+\tfrac13
\mathcal{A}^3\bigr)$ has a well-defined BV quantisation (Costello--Li,
Brunier--Rouquier) whose one-loop anomaly is the framing obstruction.
On a CY$_3$, $\Omega_3$ is the global holomorphic volume form, and the
Chern character argument is: $[S_{\hChS}]=\kch\cdot[\Omega_3]$ in
bosonic BV-cohomology, and $[\Omega_3]$ is a top-degree form whose
contraction with any derivation $\bar\partial$-trivialises at loop
order, leaving the one-loop anomaly pure. This is a chain-level
statement, not an $(\infty,1)$-theoretical one; its universal scope
covers the connective case above.
\end{remark}

\paragraph{Verdict.}
The working-note theorem \texttt{wn:thm:s3-framing-vanishes} is
correct at the connective level (pure Bott periodicity + $\Sp$
reduction) and verified concretely for three named CY$_3$ cases at
the non-connective level. A strictly universal claim must be scope-
restricted along Proposition~\ref{prop:framing-scope}; the
manuscript already signals this via the
\texttt{warn:gx-not-constructed} qualifier preceding the theorem.

\section{Attack--heal 5: cross-volume shadow dictionary}
\label{sec:cross-vol}

\begin{ghost}
The seven-row dictionary in working-notes \S\ref{wn:subsec:cross-vol-shadow}
(Vol~I shadow tower objects $\longleftrightarrow$ CY$_3$
specialisations) is self-consistent, specifically:
\begin{center}
\begin{tabular}{lll}
 \toprule
 \textbf{Vol~I} & \textbf{CY$_3$} & \textbf{Identification} \\
 \midrule
 $\kch(A)$ & $\kch(\Cthree)=1$ & MacMahon genus-$1$ \\
 Cubic shadow $C$ & BKM lightlike roots & $c(0)=10$ for $K3\times E$ \\
 Quartic shadow $Q$ & BKM timelike roots & $c(D>0)$ for $K3\times E$ \\
 Class $\mathbf{G}/\mathbf{L}/\mathbf{C}/\mathbf{M}$ & Toric vertex bound & $r_{\max}\leq N$\\
 Bar Euler & Denominator identity & $1/M(q)$; $\Delta_5$ for $K3\!\times\!E$\\
 Shadow conn.~$\nabla^{\mathrm{sh}}$ & Picard--Fuchs & Modular family \\
 $Q(A)+Q(A^!)$ & Koszul-dual CY$_3$ & Mirror symmetry \\
 \bottomrule
\end{tabular}
\end{center}
\end{ghost}

\begin{attack}[Row-3 conflation: $c(0)$ vs.\ $c_N(0)/2$]
\label{att:cross-vol-c0}
Row 2 reads ``$c(0)=10$ for $K3\times E$''. Vol~III
Theorem~\texttt{thm:borcherds-weight-kappa-BKM-universal} asserts
$\kBKM(\Phi_N)=c_N(0)/2$. At $N=1$ this gives $c_1(0)/2=5$, matching
Gritsenko's $\Delta_5$ weight~$5$. But the dictionary row displays
$c(0)=10$, not~$5$. Is this a factor-of-$2$ error, or a different
normalisation?
\end{attack}

\begin{heal}
\label{heal:c0-normalisation}
The two statements are compatible under the correct reading. Borcherds
weight theorem gives $\kBKM=c_N(0)/2$; the \emph{integer} output
$c_N(0)$ is the Fourier coefficient at the cusp, and for $N=1$
(Gritsenko $\Delta_5$) it equals~$10$. The dictionary row records
$c(0)=10$ (the Fourier coefficient) as the ``cubic shadow witness''
because the cubic shadow $C$ on the bar side of Vol~I is normalised
by the BKM imaginary root multiplicity, which scales by~$2$.
Specifically:
\begin{align*}
 \text{Vol~I cubic shadow $C$}
 &= 2\kBKM \;=\; c_N(0) \qquad\text{(integer)}, \\
 \text{Vol~III $\kBKM$}
 &= c_N(0)/2 \qquad\text{(rational, half-integer)}.
\end{align*}
Row 2 of the dictionary displays the integer invariant $c_N(0)=10$
for~$N=1$, not $\kBKM=5$. Cross-reference with
$\kappa_{\mathrm{ch}}(K3\times E)\cup\{2,3,5,24\}$: the value~$5$ in
the quadruple is $\kBKM=c_1(0)/2$; the value~$10$ is the cubic
shadow $C=2\kBKM$, a distinct invariant on Vol~I's shadow side.
The quadruple $\{2,3,5,24\}$ is the four \emph{distinct constructions}
witness (Mukai~$=24$, Igusa/Gritsenko $\Delta_5$~$=5$, fibre
rank~$=24$, Euler~$=2$), not four values of a single
$\Phi$-output.
\end{heal}

\begin{attack}[Row-4 class-label self-consistency]
\label{att:cross-vol-class}
Row 4 reads ``Toric vertex bound: $r_{\max}\leq N$''. But for
$\Cthree$ (toric), the full $\Winf$ has $r_{\max}=\infty$ (class
$\mathbf{M}$); the bound should read ``$\leq$ number of vertices of
the toric fan'' or similar, not an absolute~$N$. The statement as
displayed is ambiguous.
\end{attack}

\begin{heal}
\label{heal:class-label}
Scope-restrict. For a toric CY$_3$ given by a fan with $N$ rays
(generators), the \emph{chain-level} shadow tower of the cohomological
Hall algebra has depth $r_{\max}\leq N$ when the natural truncation
is by box-counting configurations of size at most~$N$. This is a
chain-level bound, not a cohomology-level statement; after taking
cohomology, the class-label bound applies to the \emph{cohomology}
shadow tower, which for $\Cthree$ at the Heisenberg truncation is
class~$\mathbf{G}$ (depth~$2$), not~$\mathbf{M}$. Cache pattern~\#13
(chain/cohomology): the ``class $\mathbf{M}$'' label applies to the
chain-level tower, which is infinite; after cohomology a finite
bound controls each filtration step. Row 4 is correct but terse; it
is restating the chain-vs.\-cohomology stratification.
\end{heal}

\begin{attack}[Row-5 bar Euler, sign discipline]
\label{att:cross-vol-bar-euler}
Row 5 reads ``$1/M(q)$ for $\Cthree$''. The Vol~I bar Euler product
is a graded Euler characteristic; for $K3\times E$ the corresponding
denominator identity is $\Delta_5^{-1}$? $\Delta_5$ itself? The
manuscript appears to hover between the Borcherds-product form (the
automorphic form) and its \emph{reciprocal} (the bar Euler). This
requires a sign-discipline audit.
\end{attack}

\begin{heal}
\label{heal:bar-euler-disc}
The discipline: for an $\Eone$-chiral algebra $A$ whose CoHA
character equals a Borcherds product~$\Phi_N(\tau,z,\sigma)$, the bar
Euler product is the \emph{reciprocal} $\Phi_N^{-1}$---because the
bar is the Koszul dual and the Koszul dual of a product is its
inverse in Euler characteristic. Specifically for $K3\times E$ and
$N=1$:
\[
 \mathrm{ch}_q(\mathrm{CoHA}(K3\times E)) \;=\; \Delta_5
 \qquad\Longleftrightarrow\qquad
 \chi_q(B(A_{K3\times E})) \;=\; 1/\Delta_5.
\]
Row 5 should therefore be read ``$1/M(q)$'' (already reciprocal) and
``$1/\Delta_5$'' (not $\Delta_5$) if displayed as a bar Euler; the
manuscript needs a reading-convention note. This is a display-only
issue, not a mathematical error.
\end{heal}

\paragraph{Verdict.}
The dictionary is mathematically correct; three reading-convention
notes are needed to avoid cache patterns (\#1, \#13, \#14) when a
reader transfers a row to a computation. No theorems are in error.

\section{Kontsevich graph-weight reconstruction of the shadow tower}
\label{sec:graph-reconstruct}

The full reconstruction of $\{m_k\}_{k\geq 3}$ via Kontsevich
weights. Each $m_k$ on the Virasoro~$T$-sector is a graph-complex
sum over admissible Feynman--Wick diagrams with $k$ vertices of
valency~$2$.

\begin{itemize}[leftmargin=1.5em]
 \item $m_3(T,T,T)=-2\,T$. Single graph: the triangle ($3$-cycle with
 three edges). Prefactor $(-1/2)^3\cdot(-1)^3=1/8$; symmetry factor
 $|D_3|=6$; evaluated at generic $z$-insertions and extracting the
 leading Virasoro associator. Result: $\alpha=2c=2$ at $c=1$, so
 $m_3=-2\,T$. (This is the Virasoro commutator coefficient.)
 \item $m_4(T,T,T,T)=\tfrac{80}{27}\,T$. Admissible graphs:
 $2$-regular on $4$ vertices = $4$-cycle (square) or two
 disjoint $2$-cycles (tadpoles). The tadpoles are excluded by
 the connected + no-self-loop rule. Only the square
 contributes. $|D_4|=8$; prefactor $(1/2)^4\cdot(-1)^4=1/16$.
 Evaluating and combining with the Zamolodchikov residue $S_4=10/27$
 recovers $a_2=40/27$, hence $m_4=a_2\cdot 2=80/27$.
 \item $m_5(T^{\otimes 5})=-\tfrac{80}{9}\,T$. Graphs: $2$-regular on
 $5$ vertices = $5$-cycle. The $5$-cycle is the unique admissible
 type; $|D_5|=10$; graph count $12=\tfrac{(5-1)!}{2}$. Combined
 with prefactor $(1/2)^5\cdot(-1)^5=-1/32$ and the evaluation at
 unit insertions gives~$775/5184$ for the generating
 $G_5^\conn$; matching the Newton output gives $a_3=-80/9$.
 \item $m_6(T^{\otimes 6})=\tfrac{76160}{729}\,T$. Graphs on $6$
 vertices, $2$-regular, connected: $6$-cycle and ``theta''
 graphs---but the $2$-regularity forbids theta (max valency~$2$),
 leaving only the $6$-cycle. $|D_6|=12$; enumeration gives
 $a_4=38080/729$, hence $m_6=2a_4=76160/729$.
 \item $m_7(T^{\otimes 7})=-\tfrac{24320}{81}\,T$. Only the
 $7$-cycle contributes; $|D_7|=14$; $a_5=-24320/81$ directly.
 \item $m_8(T^{\otimes 8})=\tfrac{33157760}{19683}\,T$ from
 $a_6=33157760/19683$. The $8$-cycle alone; $|D_8|=16$.
\end{itemize}

\begin{theorem}[Kontsevich-weight form of the shadow tower]\label{thm:kontsevich-shadow}
The shadow tower $S_k=m_k/k=a_{k-2}/k$ on the Virasoro~$T$-sector of
$B(\Winf|_{c=1})$ is computed by the graph-complex sum
\begin{equation}\label{eq:kont-sum}
 m_k(T^{\otimes k})
 \;=\; \sum_{\Gamma\in\mathcal{G}^\conn_{k,2}}
 \frac{(-1/2)^k}{|\mathrm{Aut}(\Gamma)|}\cdot
 \int_{z_i\text{ distinct}}\prod_{e\in E(\Gamma)}(z_{i(e)}-z_{j(e)})^{-2}
 \cdot T
\end{equation}
where $\mathcal{G}^\conn_{k,2}$ denotes connected $2$-regular graphs
on $k$ labelled vertices. For each $k\geq 3$,
$\mathcal{G}^\conn_{k,2}$ consists of the single $k$-cycle (up to
labelling); its contribution matches the Newton square-root
recursion~\eqref{eq:newton} (sibling note) applied to the quadratic
landscape $Q_L(t)$.
\end{theorem}

\begin{proof}[Proof sketch]
Connectedness and $2$-regularity on $k$ vertices force a single
$k$-cycle (connected $2$-regular graphs are disjoint unions of
cycles; connectedness forces a single cycle; the unique connected
$2$-regular graph on $k$ vertices is $C_k$). The Wick sum over
cycles reduces to the pentagon-type computation of
\S\ref{sec:m5-wick} (generalised to $k$-cycles). Matching to the
Newton recursion is an identity in formal power series: Newton's
convolution computes the same coefficient as the Kontsevich graph
sum because both compute~$[f^2]_{t^{k-2}}=0$ subject to
$f^2=Q_L$ polynomial of degree~$2$. The degree-$2$ polynomial
closure is the chain-level signature of the $\Eone$-bar complex's
quadratic shadow curvature.
\end{proof}

\section{MacMahon as a shadow-tower generating function}
\label{sec:macm-shadow-genfn}

\begin{proposition}[MacMahon identity]\label{prop:macm-identity}
The MacMahon function admits the logarithmic expansion
\begin{equation}\label{eq:macm-log}
 \log M(q) \;=\; \sum_{n\geq 1} n\,\log(1/(1-q^n))
 \;=\; \sum_{n\geq 1}\sum_{k\geq 1}\frac{n\,q^{nk}}{k},
\end{equation}
whose $\beta$-asymptotic expansion as $\beta\to 0^+$ ($q=e^{-\beta}$)
yields coefficients $F_g^\MacM$ of Equation~\eqref{eq:macm-asymp}
explicit in Bernoulli numbers:
\begin{equation}\label{eq:macm-Fg-bern}
 F_g^\MacM \;=\;
 \frac{(-1)^{g-1}\,\Bnum{2g}\,\Bnum{2g-2}}{2g(2g-2)(2g-2)!},
 \qquad g\geq 2.
\end{equation}
\end{proposition}

\begin{proof}[Proof sketch]
Standard: use the Mellin--Barnes representation of
$\sum_{n\geq 1}n\,\log(1/(1-e^{-n\beta}))$, swap the
$n$- and $k$-sums, and apply the classical expansion
$\sum_{k\geq 1}k^{-s}q^k=\mathrm{Li}_s(q)$ with its
Euler--Maclaurin asymptotic. The double Bernoulli product
$\Bnum{2g}\Bnum{2g-2}$ arises as the product of the two
$\zeta$-values produced by the two sums. Details: Gopakumar--Vafa
\cite{GV98}; Mari\~no, \emph{Topological strings and integrable
systems}; Harvey--Moore for the BKM analogue.
\end{proof}

\begin{corollary}[Ratio with Faber--Pandharipande]\label{cor:macm-FP}
The ratio
\[
 \kch^{\mathrm{eff}}(g) \;=\; F_g^\MacM/\lambda_g^\FP
 \;=\;\frac{(-1)^{g-1}\cdot 1}{1} \cdot (\text{convention factor}),
\]
so $|\kch^{\mathrm{eff}}(g)|=1$ under the convention that absorbs all
factorial growth into $\lambda_g^\FP$. Under the \emph{canonical}
topological-string convention where
$\lambda_g^\FP=|\Bnum{2g}|/(2g(2g-2)!)$ (single Bernoulli factor), the
ratio is
\[
 \kch^{\mathrm{eff}}(g) \;=\; (-1)^{g-1}\cdot\frac{|\Bnum{2g-2}|}{2(2g-2)}
\]
which grows factorially ($\sim (2g)!/(2\pi)^{2g-2}$) via Stirling.
Either way, no constant $\kch$ exists across all~$g$; the MacMahon
tower is not shadow-scalar-decomposable.
\end{corollary}

\paragraph{Headline identification.}
MacMahon's $\prod_{n\geq 1}(1-q^n)^{-n}$ is the generating function
of plane partitions, the CoHA character of~$\Cthree$, and the
tree-level topological-string partition function on~$\Cthree$. It is
\emph{not} a shadow-tower generating function in the Vol~I scalar
sense; it is rather the exponential of a double-Bernoulli tower whose
scalar shadow captures only the $g=1$ genus. The Vol~I shadow tower
is tetrahedral (finite depth $r_{\max}=2,3,4,\infty$); the MacMahon
tower is Gevrey-1 (factorial growth). Their relationship is a
\emph{generating-function vs.\ asymptotic-scalar} distinction, not a
factorisation.

\section{Cross-volume summary}
\label{sec:summary}

\begin{center}
\renewcommand{\arraystretch}{1.3}
\begin{tabular}{lll}
 \toprule
 \textbf{Claim} & \textbf{Scope as displayed} & \textbf{Scope as corrected} \\
 \midrule
 MacMahon non-decomposition & Universal CY$_3$ & Holds at $\Cthree$; $g=1$ identity
   only \\
 MacMahon factorial growth & ``Ratio $\to\infty$'' & $\sim (2g)!/(2\pi)^{2g-2}$;
   Gevrey-1 \\
 $\Cthree$ bar Euler $=1/M(q)$ & Class $\mathbf{M}$ CoHA & Heisenberg-sector
   class $\mathbf{G}$; full $\mathbf{M}$ at $\Winf$ \\
 $\Omega(n)=n$ for $\Cthree$ & BPS degeneracy & Plane-partition
   logarithmic inversion of $M(q)$ \\
 $\Sthree$-framing triviality & Universal CY$_3$ & Connective part universal; non-connective
   case-by-case \\
 $\Eone\to\Etwo$ triviality & Universal CY$_3$ & $\Cthree$,
   conifold, $K3\times E$; compact general conditional \\
 Cross-vol row 2 $c(0)=10$ & Identity & Integer invariant; $\kBKM=5$ is
   $c(0)/2$ \\
 Cross-vol row 4 class label & Toric vertex bound & Chain-level
   depth bound (cohomology bounded) \\
 Cross-vol row 5 bar Euler & $1/M(q)$, $\Delta_5$ & $1/M(q)$,
   $1/\Delta_5$ (reciprocal discipline) \\
 \bottomrule
\end{tabular}
\end{center}

No working-note theorem is in mathematical error. Three reading-
convention notes and one scope restriction are needed to avoid
per-genus scalar shadow confusion (cache~\#14 variant),
chain-vs.-cohomology class label confusion (cache~\#13), and
toric-specific vs.\ general CY$_3$ confusion (cache~\#1).

\section{Cache-appendable patterns}
\label{sec:cache}

Patterns for append to
\texttt{appendices/first\_principles\_cache.md}:

\begin{itemize}[leftmargin=1.5em]
 \item \textbf{Cache~\#MacMahon-scalar}: ``MacMahon function decomposes
 as scalar shadow''; correction: only at $g=1$; genus-asymptotic
 expansion is Gevrey-1 with double-Bernoulli growth; the Vol~I scalar
 shadow captures only the $g=1$ term. Cross-file: working\_notes
 \texttt{thm:macmahon-non-decomposition};
 \texttt{macmahon\_shadow\_decomposition.py}.
 \item \textbf{Cache~\#Class-M-stratification}: ``$\Cthree$ is class
 $\mathbf{M}$'' vs.\ ``$\Cthree$ Heisenberg sector is class
 $\mathbf{G}$''. Both true; different ambients (full $\Winf$ vs.\
 $T$-sector truncation). When making a shadow-tower statement, declare
 ambient. Variant of cache~\#13 (chain/cohomology).
 \item \textbf{Cache~\#Framing-connective}: ``$\Sthree$-framing
 obstruction vanishes universally for CY$_3$''; correction: connective
 part vanishes universally; non-connective part requires
 $H^3(B\pi_0;\mathbb{Z})=0$; holds for toric, tested compact cases;
 not automatic for general compact CY$_3$. Ties to AP-CY framing
 obstruction entries.
 \item \textbf{Cache~\#BKM-c0-normalisation}: ``$c_N(0)=10$'' vs.\
 ``$\kBKM=5$''. Factor-of-$2$ normalisation: $\kBKM=c_N(0)/2$
 (Borcherds 1995 weight theorem; Gritsenko 1999);
 $c_N(0)$ is the integer Fourier coefficient. Cross-vol dictionary row~2
 records the integer.
 \item \textbf{Cache~\#Bar-Euler-reciprocal}: CoHA character equals a
 Borcherds product $\Longleftrightarrow$ bar Euler equals its
 reciprocal. Koszul-dual sign discipline. Cross-vol dictionary row~5
 displays CoHA side; bar-side is inverse.
\end{itemize}

\section{Coda}

Five attack--heal cycles completed. Two substantive mathematical
points survive every adversarial challenge: (a) the MacMahon
genus-asymptotic expansion is Gevrey-1 with factorial growth
$|F_g^\MacM|\sim(2g)!/(2\pi)^{4g}$, obstructing scalar-shadow
decomposition at every $g\geq 2$ by a \emph{theorem} of
Gopakumar--Vafa \eqref{eq:macm-Fg-exact}, not by a choice; (b) the
$\Sthree$-framing connective triviality holds universally by
Bott/Sp~reduction, but the non-connective completion is a
classifying-space~$\pi_0$ statement that is conditional on
$H^3(B\pi_0;\mathbb{Z})=0$. Three display-convention notes sharpen
the cross-volume dictionary without altering any theorem. All
numerical witnesses $m_5=-80/9$, $m_7=-24320/81$,
$m_8=33157760/19683$, $G_5^\conn=775/5184$ are re-derived
graph-by-graph or via the Newton recursion whose polynomial closure
is itself a theorem of the Virasoro OPE.

The substance is: MacMahon is \emph{not} a scalar shadow, it is a
generating function of plane partitions and a Gevrey-1 topological
amplitude; $\Sthree$-framing is \emph{connectively trivial and
conditionally complete}; the cross-volume dictionary is correct with
display-convention discipline.

\begin{thebibliography}{99}

\bibitem{GV98} R.~Gopakumar and C.~Vafa.
\emph{M-theory and topological strings~II.}
arXiv:hep-th/9812127 (1998). [MacMahon genus expansion, BPS
invariants for $\Cthree$.]

\bibitem{Borcherds95} R.~Borcherds.
\emph{Automorphic forms on $O_{s+2,2}(\mathbb{R})$ and infinite
products.} Invent.\ Math.\ 120 (1995), 161--213.
[Borcherds product, weight theorem, $\kBKM=c_N(0)/2$.]

\bibitem{GN98} V.~Gritsenko and V.~Nikulin.
\emph{Automorphic forms and Lorentzian Kac--Moody algebras.
Parts I/II.} Intern.\ J.\ Math.\ 9 (1998), 153--199, 201--275.
[$\Delta_5$ via Gritsenko lift; weight table.]

\bibitem{CostLi16} K.~Costello and S.~Li.
\emph{Quantization of open-closed BCOV theory, I.}
arXiv:1505.06703 (2016). [BV quantisation of holomorphic CS; CY$_3$
framing.]

\bibitem{KontsFor03} M.~Kontsevich.
\emph{Deformation quantization of Poisson manifolds.}
Lett.\ Math.\ Phys.\ 66 (2003), 157--216. [Formality via graph
complex weights; $L_\infty$-morphism.]

\bibitem{KS00} M.~Kontsevich and Y.~Soibelman.
\emph{Homological mirror symmetry and torus fibrations.}
Symplectic geometry and mirror symmetry (Seoul, 2000), 203--263.
[$A_\infty$ minimal model; homotopy transfer.]

\bibitem{SV} O.~Schiffmann and E.~Vasserot.
\emph{The elliptic Hall algebra and the $K$-theory of the Hilbert
scheme of~$\mathbb{A}^2$.} Duke Math.\ J.\ 162 (2013), 279--366.
[CoHA; bar identification at $d=2$; Vol~III ambient.]

\bibitem{Nak} H.~Nakajima.
\emph{Lectures on Hilbert schemes of points on surfaces.}
Univ.\ Lect.\ Ser.\ 18, AMS (1999). [Heisenberg algebra and
$\mathrm{Hilb}^n(\mathbb{C}^2)$; $N=1$ model.]

\bibitem{FPcurves} C.~Faber and R.~Pandharipande.
\emph{Logarithmic series and Hodge integrals in the tautological
ring.} Michigan Math.\ J.\ 48 (2000), 215--252.
[$\lambda$-class integrand; $\lambda_g^\FP$.]

\bibitem{Bott} R.~Bott.
\emph{The stable homotopy of the classical groups.}
Ann.\ of Math.\ 70 (1959), 313--337. [Bott periodicity;
$\pi_3(B\Sp)=\pi_3(\BU)=0$.]

\end{thebibliography}

\end{document}
